\documentclass[../../main.tex]{subfiles}
\begin{document}

{\bf [解答]}

% TODO: 概要図

条件(1)から，
\begin{align}
  a+b=3 \label{eq:1}
\end{align}
である．次に条件(2)を考える．$x=0$のときはこれは満たされるから，以下$x>0$で考える．このとき，
\begin{align}
  \forall x>0,f(x)\ge0 \nonumber\\
  \Longleftrightarrow \forall x>0,x^2+ax+b\ge0 \label{eq:2}
\end{align}
だから\cref{eq:2}について考えればよい．そこでこの不等式の左辺を$g(x)$，
方程式$g(x)=0$の判別式を$D$とおけば，\cref{eq:2}は
\begin{align}
  &D\le0 \lor
  \begin{cases}
    D>0  \\
    g(0)\ge0 \\
    \dfrac{-a}{2}< 0
  \end{cases} \\
  \Longleftrightarrow
  &a^2-4b\le 0 \lor
  \begin{cases}
    a^2-4b>0  \\
    b\ge0 \\
    a>0
  \end{cases} \\
  \Longleftrightarrow
  &-6\le a\le2 \lor
  \begin{cases}
    a<-6,2<a  \\
    0<a\le3
  \end{cases} \quad (\because\text{\cref{eq:1}}) \\
  \Longleftrightarrow
  &-6\le a\le 3 \label{eq:3}
\end{align}
となる．このもとで$h(a)=\int_0^1f(x)dx$の最大・最小を考える．
\begin{align}
  h(a)&=\left[ \dfrac{1}{4}x^4+\dfrac{a}{3}x^3+\dfrac{b}{2}x^2\right]_0^1 \\
  &=\dfrac{-1}{6}a+\dfrac{7}{4}\tag{$\because\cref{eq:1}$}
\end{align}
となり，$h(a)$は$a$の単調減少関数である．故に\cref{eq:1,eq:3}より，$(a,b)=(3,0)$で最小，$(a,b)=(-6,9)$で最大となる．$\cdots$(答)
\newpage
\end{document}
